\documentclass[11pt]{article}
\usepackage[margin=1in]{geometry}
\usepackage{amsmath,amssymb,amsthm,hyperref,enumitem}
\newtheorem{theorem}{Theorem}
\newcommand{\Si}{\operatorname{Si}}
\newcommand{\Ci}{\operatorname{Ci}}
\title{Both endpoint asymptotics of the first zero of $E_{a,a}(-t^a)$}
\author{Principia Math In Progress}
\date{21 September 2026}
\begin{document}
\maketitle
\noindent\textbf{Task 1: MathDB 361624.}
Put
\[
 H(a,t)=E_{a,a}(-t^a)=\sum_{k\geq0}
       \frac{(-1)^kt^{ak}}{\Gamma(a(k+1))},
 \qquad Z(a)=\min\{t>0:H(a,t)=0\}.
\]
Conjecture 1(b) of Chaudhary, Diethelm and Hashemishahraki
\cite{CDH} proposes $Z(a)=(c\log(a-1)+d)(1+o(1))$ as
$a\downarrow1$, with $c\approx-0.81$ and $d\approx2.25$.
Part (c) proposes $Z(a)=\pi+c(2-a)+o(2-a)$ as $a\uparrow2$,
again with $c\approx-0.81$. The authors explicitly ask whether these
two coefficients are equal. The following theorem answers this question
and corrects the first coefficient.

\begin{theorem}\label{main}
The first positive zero exists for all $a>1$ sufficiently close to
either endpoint $1$ or $2$. With $L=\log(1/\varepsilon)$,
\begin{align*}
 \boxed{Z(1+\varepsilon)=L+2\log L+o(1)}
       &\qquad(\varepsilon\downarrow0),\\
 \boxed{Z(2-\delta)=\pi+c_*\delta+O(\delta^2),\quad
 c_*={\pi\Ci(\pi)-\Si(\pi)\over2}}
       &\qquad(\delta\downarrow0).
\end{align*}
Here $-811/1000<c_*<-809/1000$, so the two leading coefficients
are unequal: the first is $-1$, and the second is
$-0.8102512403460584\ldots$.
\end{theorem}

We use $\Si(x)=\int_0^x\sin u/u\,du$ and
$\Ci(x)=\gamma+\log x+\int_0^x(\cos u-1)/u\,du$ for $x>0$.
The theorem refutes the numerical leading coefficient in part (b)
and hence the conjunction recorded in MathDB. It also supplies the
exact coefficient in part (c). The global positivity bound
$H(a,t)>0$ for $0\leq t\leq2.93785$ in part (a) remains unresolved.
No contradiction rests on the constant $d$, which is invisible at
the source's multiplicative $1+o(1)$ precision.

\section*{An exact representation near order one}
Write $a=1+\varepsilon$, $0<\varepsilon<1/4$,
$\theta=\pi\varepsilon$, and
\[
 K(r)={\sin\theta\over\pi}
 {r^{a-1}\over (r^a-1)^2+2r^a(1-\cos\theta)}\quad(r>0).
\]
For $f_a(t)=E_{a,1}(-t^a)$ the classical residue formula is
\begin{equation}\label{f}
 f_a(t)={2\over a}e^{t\cos(\pi/a)}\cos(t\sin(\pi/a))
             -\int_0^\infty e^{-tr}K(r)\,dr.
\end{equation}
Indeed, termwise Laplace transformation for real $s>1$ gives
$s^{a-1}/(s^a+1)$. On its principal branch the poles are
$e^{\pm i\pi/a}$, with residues in the inverse transform
$a^{-1}e^{t e^{\pm i\pi/a}}$. The imaginary part of the transform
on the upper side of the negative real cut is
\[
 \operatorname{Im}{r^{a-1}e^{i\theta}\over1-r^ae^{i\theta}}
       =\pi K(r).
\]
Deforming the inversion contour gives the negative integral in
\eqref{f}. The small-circle contribution is $O(\rho^a)$ and the
large arcs vanish for $t>0$, using the transform's $O(1/s)$ decay
and the exponential on the left arcs. The cut integral converges
absolutely. These representation methods are classical; see
\cite{Hanneken,Paris}.

Differentiating the defining series gives
$-f_a'(t)=t^{a-1}H(a,t)$. Differentiating the exact integral in
\eqref{f} is justified for $t>0$ by exponential integrability.
Thus, with $g_a(t)=t^{a-1}H(a,t)$,
\begin{equation}\label{g}
 g_a(t)=-{2\over a}e^{t\cos(\pi/a)}
       \cos\bigl(t\sin(\pi/a)+\pi/a\bigr)
             -\int_0^\infty r e^{-tr}K(r)\,dr.
\end{equation}
This is differentiation of an exact identity, not of an asymptotic
expansion. The unweighted mass, by $u=r^a$, is
\begin{equation}\label{mass}
 \int_0^\infty K(r)\,dr
 ={\sin\theta\over\pi a}\int_0^\infty
 {du\over(u-\cos\theta)^2+\sin^2\theta}
 ={\pi-\theta\over\pi a}={2\over a}-1.
\end{equation}

\section*{Uniform estimates and the first zero near order one}
We prove, uniformly for $L/4\leq t\leq3L/2$,
\begin{equation}\label{uniform}
 g_a(t)=e^{-t}(1+o(1))
          -\varepsilon\Gamma(a+1)t^{-a-1}(1+o(1)),
\end{equation}
with separate relative errors for the two displayed terms.
All constants below are uniform for sufficiently small $\varepsilon$.
Write $D(r)=1-2r^a\cos\theta+r^{2a}$ and $h=t^{-2}$.
For $r\leq1/2$, $D\geq1/4$ and $D^{-1}=1+O(r^a)$, so
\begin{align*}
 \int_0^{1/2}r e^{-tr}K(r)\,dr
 &= {\sin\theta\over\pi}\left\{
 \Gamma(a+1)t^{-a-1}+O(t^{-2a-1})
       +O\bigl(t^{-a-1}(1+t)^2 e^{-t/2}\bigr)\right\}\\
 &=\varepsilon\Gamma(a+1)t^{-a-1}(1+o(1)).
\end{align*}
The algebraic error follows by extending its gamma integral to
infinity. The omitted main-integral tail follows by $u=tr$ and
$u^a\leq1+u^2$, since $1<a<5/4$.

Outside $[1/2,3/2]$ the unweighted $K$ mass is $O(\varepsilon)$:
on the lower interval use $D\geq1/4$; on the upper interval use
$D\geq(r^a-1)^2\geq r^{2a}/9$, whence
$K(r)\leq C\varepsilon r^{-a-1}$. On $[1/2,3/2]$ the mean value
theorem yields $|r^a-1|\geq C_0|r-1|$, so
$K(r)\leq C\varepsilon/(r-1)^2$ for $r\ne1$. Equation \eqref{mass}
therefore implies
\[
 \int_{1-h}^{1+h}K(r)\,dr
       =1+O(\varepsilon)+O(\varepsilon/h)=1+o(1).
\]
Here $\varepsilon t^2\to0$. In this central interval,
$r e^{-tr}=e^{-t}(1+O(h+th))$, and consequently its contribution
to the integral in \eqref{g} is $e^{-t}(1+o(1))$.
The rest of $[1/2,3/2]$ contributes at most
$C\varepsilon t^4e^{-t/2}=o(e^{-t})$, since its ratio to
$e^{-t}$ is at most $C L^4\varepsilon^{1/4}$.
For the remaining upper tail we retain the exponential weight:
\[
 \int_{3/2}^\infty r e^{-tr}K(r)\,dr
 \leq C\varepsilon\int_{3/2}^\infty r^{-a}e^{-tr}\,dr
 \leq C\varepsilon e^{-3t/2}/t=o(e^{-t}).
\]
Finally $\cos(\pi/a)=-1+O(\varepsilon^2)$ and
$\sin(\pi/a)=O(\varepsilon)$. The residue term in \eqref{g} is
$2e^{-t}(1+o(1))$ because $\varepsilon t\to0$ uniformly.
Subtracting the central and small-$r$ contributions proves
\eqref{uniform}.

We next exclude earlier zeros. For $1\leq b\leq a$ and $k\geq1$,
\[
 \partial_b {t^{bk}\over\Gamma(b(k+1))}
 ={t^{bk}\over\Gamma(b(k+1))}
        \bigl(k\log t-(k+1)\psi(b(k+1))\bigr).
\]
We have $\Gamma(b(k+1))\geq k!$ and
$0<\psi(b(k+1))\leq\log(b(k+1))\leq2(k+1)$.
The $k=0$ derivative is uniformly bounded. For $1\leq t\leq L/4$,
eventually $t^{b-1}\leq2$, and summing the factorial majorant gives
$|\partial_b H(b,t)|\leq C(1+t)^3e^{2t}$.
As $H(1,t)=e^{-t}$, integration over $b$ yields
\[
 { |H(a,t)-e^{-t}|\over e^{-t}}
 \leq C\varepsilon(1+L)^3e^{3L/4}=o(1).
\]
For $0<t\leq1$, use $k t^{bk}|\log t|\leq1/(b\mathrm e)$;
the derivative series is uniformly bounded by a summable multiple
of $(1+(k+1)^2)/k!$. Hence $H(a,t)=e^{-t}+O(\varepsilon)$
uniformly there. Also $H(a,0)=1/\Gamma(a)>0$.
These bounds justify the differentiations and show positivity on
$[0,L/4]$ for small $\varepsilon$.

Fix $0<\eta<1/2$. The ratio of the second leading term in
\eqref{uniform} to the first is
\[
 Q_\varepsilon(t)=\varepsilon\Gamma(a+1)e^t t^{-a-1}.
\]
It is increasing for $t>a+1$. At $(1-\eta)L$ it tends to zero,
and at $(1+\eta)L$ it tends to infinity. Equation \eqref{uniform}
therefore gives positivity throughout $[L/4,(1-\eta)L]$ and a
negative value at $(1+\eta)L$. Since $g$ and $H$ have the same
sign for $t>0$, continuity and earlier positivity show that a first
positive zero exists and $Z(a)/L\to1$.
At this zero the two relative errors in \eqref{uniform} tend to zero.
Taking logarithms gives
\[
 Z(a)=L+(a+1)\log Z(a)-\log\Gamma(a+1)+o(1)
      =L+2\log L+o(1).
\]
The last step uses $\log Z(a)=\log L+o(1)$,
$\varepsilon\log Z(a)=o(1)$, and $\log\Gamma(a+1)=o(1)$.

\section*{The coefficient near order two}
Set $F(a,t)=tH(a,t)$. Then $F(2,t)=\sin t$ and
$F_t(2,\pi)=-1$. The series and its first two derivatives converge
locally uniformly near $(2,\pi)$: for $a\geq3/2$ and $k\geq1$,
$\Gamma(a(k+1))\geq k!$, and derivatives contribute only
polynomial and logarithmic factors in $k$. The $k=0$ term is smooth.
The implicit function theorem gives a $C^2$ zero branch $z(a)$,
and its coefficient in $2-a$ is $-F_a(2,\pi)$.
For small fixed $\eta>0$, $F_t(a,t)<-1/4$ near
$[\pi-\eta,\pi+\eta]$ when $a$ is close to $2$.
On $[0,\pi-\eta]$, $H(a,t)$ converges uniformly to
$\sin t/t$ (value $1$ at zero), which has a positive minimum.
Thus the branch is the first positive zero.

To evaluate its derivative, introduce
$F_2(a,t)=tE_{a,2}(-t^a)$ and $w(t)=\partial_a F_2(2,t)$.
Termwise differentiation and $\psi(m+1)=H_m-\gamma$ give
\[
 w(t)=\sum_{k\geq1}{(-1)^k k t^{2k+1}\over(2k+1)!}
             (\log t+\gamma-H_{2k+1})
 =-{1\over2}\int_0^t(\log(t-u)+\gamma)u\sin u\,du.
\]
The integral identity follows by expanding $u\sin u$ and using
$\int_0^t u^m\log(t-u)\,du=
t^{m+1}(\log t-H_{m+1})/(m+1)$.
With $v=\pi-u$, two integrations by parts yield
\[
 -w(\pi)={1\over2}\int_0^\pi(\log v+\gamma)(\pi-v)\sin v\,dv
         ={\pi\Ci(\pi)+\Si(\pi)\over2}.
\]
For explicit primitives, use
$-\log v\cos v+\Ci(v)$ for $\log v\sin v$ and
$-v\log v\cos v+(\log v+1)\sin v-\Si(v)$ for
$v\log v\sin v$. Their limits at zero are $\gamma$ and $0$.

The derivative of the second parameter in $F$ contributes one extra
digamma factor. Thus the required coefficient is
\begin{align*}
 c_*&=-w(\pi)+\sum_{k\geq0}
       {(-1)^k\pi^{2k+1}\psi(2k+2)\over(2k+1)!}\\
 &=-w(\pi)+\int_0^1{\sin\pi-\sin(\pi u)\over1-u}\,du
 =-w(\pi)-\Si(\pi).
\end{align*}
Here the Euler-constant part vanishes, and
$H_{2k+1}=\int_0^1(1-u^{2k+1})/(1-u)\,du$ justifies the
middle equality by absolute convergence of the integrated series.
This proves the closed formula; Taylor's theorem gives the stated
$O(\delta^2)$ remainder.

\section*{An exact rational enclosure}
We specify a finite rational certificate for $c_*\ne-1$.
Let $p_-=31415926535/10^{10}$ and $p_+=31415926536/10^{10}$.
Machin's identity $\pi=16\arctan(1/5)-4\arctan(1/239)$,
with 12 alternating terms and the next-term bound for each arctangent,
proves $p_-<\pi<p_+$.
For rational $x>1$ put $y=(x-1)/(x+1)$ and
\[
 A(x)=2\sum_{j=0}^{49}{y^{2j+1}\over2j+1},\qquad
 B(x)=A(x)+{2y^{101}\over101(1-y^2)}.
\]
Then $A(x)<\log x<B(x)$. With $N=2048=2^{11}$,
\[
 H_N-11B(2)-1/N<\gamma<H_N-11A(2),\qquad
 A(p_-)<\log\pi<B(p_+).
\]
The gamma bound follows from
$0<H_N-\log N-\gamma<1/N$, obtained by summing
$0<\log(1+1/j)-1/(j+1)<1/j-1/(j+1)$ for $j\geq N$.
Use the alternating series
\begin{align*}
 \Ci(\pi)&=\gamma+\log\pi+
       \sum_{k=1}^{10}{(-1)^k\pi^{2k}\over2k(2k)!}+R_C,
 &-{p_+^{22}\over22\cdot22!}<R_C<0,\\
 \Si(\pi)&=\sum_{k=0}^{10}
       {(-1)^k\pi^{2k+1}\over(2k+1)(2k+1)!}+R_S,
 &-{p_+^{23}\over23\cdot23!}<R_S<0.
\end{align*}
The term magnitudes decrease, as follows from their ratios and
$\pi^2<10$. Evaluate positive powers in $[p_-^j,p_+^j]$, add
intervals with their displayed signs, and use all four endpoint
products when multiplying intervals. The exact resulting rational
endpoints $c_-,c_+$, stored by the verifier, satisfy
\[
 -811/1000<c_-<c_*<c_+<-809/1000.
\]
Together with the first asymptotic, this proves the coefficients
unequal and completes Theorem~\ref{main}.
\hfill$\square$

\section*{Verification, attribution, and scope}
The representation and perturbation methods are classical. The
uniform weighted-kernel estimates above and the first-zero arguments
are included explicitly. Companion MathDB 361626 treats second
parameter $1$ near order $1$; MathDB 361625 treats fixed second
parameter $2$ near order $2$. Here the second parameter varies with
the order, changing both coefficients. These are distinct source
conjectures. Earlier first-zero existence and monotonicity results
include \cite{EL}; no novelty for those general facts is claimed.

The verifier checks rational constants, the coefficient enclosure,
and symbolic identities exactly. Supplemental 110-digit calculations
compare the exact integral with the series and examine both endpoint
limits. These are not interval certificates or a substitute for
first-zero minimality. The complete current primary and author-posted
target-final text were read; limitations of the bounded priority
audit are recorded separately. Review is by the authoring agent only.

\begin{enumerate}[label=\textbf{Task \arabic*:},start=2,leftmargin=*]
\item Determine the next term after $L+2\log L$, with an error bound.
\item Prove or refute the global positivity bound in Conjecture 1(a).
\item Determine the quadratic coefficient of $(2-a)^2$ in $Z(a)$.
\item Prove uniqueness and locate the global minimum of $Z(a)$
on $(1,2]$ with rigorous bounds.
\end{enumerate}
\small
\begin{thebibliography}{9}
\bibitem{CDH} R. Chaudhary, K. Diethelm and S. Hashemishahraki,
On the separation of solutions to fractional differential equations
of order $\alpha\in(1,2)$, \emph{Applied Numerical Mathematics}
\textbf{203} (2024), 84--96,
\href{https://doi.org/10.1016/j.apnum.2024.05.020}{doi:10.1016/j.apnum.2024.05.020}.
\href{https://arxiv.org/abs/2401.14771}{Current preprint: arXiv:2401.14771v3}.
Conjecture 1 in both versions.
\bibitem{Hanneken} J. W. Hanneken, B. N. N. Achar and D. M. Vaught,
An alpha-beta phase diagram representation of the zeros and properties
of the Mittag-Leffler function, \emph{Advances in Mathematical Physics}
(2013), Article 421685,
\href{https://doi.org/10.1155/2013/421685}{doi:10.1155/2013/421685}.
\bibitem{Paris} R. B. Paris, Exponential asymptotics of the
Mittag--Leffler function, \emph{Proceedings of the Royal Society A}
\textbf{458} (2002), 3041--3052,
\href{https://doi.org/10.1098/rspa.2002.0977}{doi:10.1098/rspa.2002.0977}.
\bibitem{EL} P. W. Eloe and Y. Li, Existence, comparison, and
monotonicity of principal eigenvalues of fractional Sturm--Liouville
problems involving Riemann--Liouville derivatives,
\emph{Fractional Calculus and Applied Analysis} \textbf{29} (2026),
1--25, \href{https://doi.org/10.1007/s13540-026-00481-1}{doi:10.1007/s13540-026-00481-1}.
\end{thebibliography}
\end{document}
